\section{Results}
\label{sec:sn:results}

\subsection{Exploration by elimination}
HCA-DSE traverses architectures as a tree of partial designs and applies three tests of
increasing cost at every node (Fig.~\ref{fig:ce:concept}). A partial design is discarded
when it is structurally ill-formed, when an optimistic bound on any constrained quantity
already violates its limit, or when an optimistic bound on the whole objective vector is
dominated by an architecture that has already been evaluated. Only designs that survive
to a leaf are handed to the expensive evaluator. The bounds are admissible by
construction: each additive term of an objective is relaxed independently over the
domains still reachable, and the waiting term becomes exact once the determinants of tier
utilisation have been assigned (Methods). Safety of the two bound-based rules and
preservation of the exact front follow from this property; the statements and proofs are
in Methods and Supplementary Note~1.

\subsection{Expensive evaluations become the scarce resource that is actually spent}
Across four design-space scales and three workload classes, the funnel removes the design
space before it is evaluated rather than after (Fig.~\ref{fig:ce:funnel}). On a space of
$51\,840$ candidates the method performs $\SthreeEvalRange$ full evaluations, discarding
$99.90$--$99.92\%$ of the raw space beforehand; on $345\,600$ candidates it performs
$\SfourEvalRange$, a reduction of $99.98$--$99.99\%$. The composition differs by
workload: for data-intensive sensing the feasibility bounds remove the largest share,
whereas for telemetry dominance pruning does most of the work, reflecting which
constraint binds in each regime.

The reduction is exact, not approximate. For every scale at which exhaustive enumeration
is tractable and for every workload, the set of objective vectors returned by HCA-DSE is
identical to the exhaustive Pareto front, with recall and precision of $1.0$
(Fig.~\ref{fig:ce:pareto}). Admissibility of the bounds was verified exhaustively --- all
partial designs against all their completions --- and for every pruned partial design all
completions were enumerated and checked against the claim of the rule that discarded it.
An independent exact solver built on a finite table and an SMT backend recovers the same
front in all six small and medium cases, at the price of evaluating the entire
structurally valid space first: $1080$ and $3888$ evaluator calls respectively
(Supplementary Table~S5).

\subsection{The saving converts into time exactly when evaluation is expensive}
Whether fewer evaluations translate into less time depends on what an evaluation costs.
We measured the overhead of the exploration itself and the cost of one design-point
evaluation for two evaluators, and projected the resulting speed-up
(Fig.~\ref{fig:ce:crossover}). With a closed-form analytical evaluator, whose cost is
comparable to that of the bounds, the measured speed-up over exhaustive enumeration is
only $1.32$--$3.15\times$ despite the $99.9\%$ reduction in evaluations: in that regime
pruning is not the right tool, and we report it as such. With a discrete-event evaluator
measured at $\DesCostMs$\,ms per design point, the same reduction projects to
$86.4$--$146.7\times$ on the medium space and $638.2$--$771.1\times$ on the large one.
The break-even evaluator cost --- the price above which pruning repays its own overhead
--- is $1.25$--$2.55\,\mu$s for the medium space and at or below $0.40\,\mu$s for the
large one, on the hardware and implementation used here. These are projections computed
from measured counts and measured overheads; neither strategy was executed end to end
with the discrete-event evaluator in the inner loop, and the break-even value is a
property of this implementation rather than a constant of the method.

\subsection{Guarantees at a fraction of a stochastic search budget}
Under matched evaluation budgets, NSGA-II approaches the exact front but does not reach
it, and cannot certify what is missing (Fig.~\ref{fig:ce:hv}, Table~\ref{tab:ce:base}).
At a budget of $1000$ full evaluations on the large space its median Pareto recall is
$0.500$--$0.741$ with a hypervolume ratio of $0.946$--$0.984$, while random search
recovers at most $3.2\%$ of the front points. HCA-DSE returns the complete front with
$\SthreeEvalRange$ evaluations, between $20$ and $24$ times fewer. On the medium space
NSGA-II is closer in hypervolume ($0.999$--$1.000$) yet still misses $7$--$20\%$ of the
front points. Across the six exploratory comparisons between NSGA-II and random search at
the two largest budgets, none reaches $p<0.05$ (unadjusted, $0.057$--$0.338$; Cliff's
$\delta$ $0.080$--$0.295$), so we do not claim a separation between the two stochastic
baselines --- the comparison that matters here is between guaranteed and approximate
outcomes, not between two approximations.

\subsection{The mechanism scales because the space is never built}
Sweeping the raw space over five orders of magnitude, from $1.3\times10^{3}$ to
$8.3\times10^{7}$ candidates, the number of expanded partial designs never exceeds
$16\,780$ and the number of full evaluations stays between $18$ and $61$
(Fig.~\ref{fig:ce:scal}). The largest observed search time is $3.98$\,s with memory
tracing enabled, and peak traced allocation is $99.6$\,KiB, because the Cartesian product
is generated lazily and never materialised. Exhaustive front equality cannot be checked
at these sizes and is not claimed there.

Two internal ablations show where the reduction comes from. Removing the bound layers one
at a time on the large space leaves $34\,560$ evaluations for structural pruning alone,
$7025$--$12\,700$ with feasibility bounds, $322$--$415$ with dominance bounds, and
$\SthreeEvalRange$ with both --- the layers are complementary, because feasibility
pruning removes exactly the sub-trees in which the dominance test would otherwise be
repeated (Supplementary Table~S6). Toggling the contention-aware term of the bound while
holding the traversal order fixed isolates its contribution: $3777$--$5139$ incomplete
designs carry a positive queueing contribution, of which $163$--$419$ are discarded only
because that term is present. All $36$ order/toggle/scale/workload combinations preserve
the exact front (Supplementary Table~S7).

\subsection{The ordering transfers to a real deployment}
To test whether an architecture ranking produced at design time survives contact with a
real system, we deployed the same architecture template as containers --- one CPU per
tier node, per-node link shaping with delay and bandwidth of the design under test, and a
wide-area profile on the cloud hop --- and measured $36$ architectures, three repeats
each, $108$ runs (Table~\ref{tab:ce:testbed}, Fig.~\ref{fig:ce:testbed}). The analytical
model was frozen before this round: it was not refitted to these observations.

Absolute prediction error is $25.4\%$ (pooled MAPE) with a median absolute error of
$6.51$\,ms. The decision-relevant quantity is different: among architecture pairs that
the model separates by more than $25\%$, the measured order agrees with the predicted
order for $88.6$--$91.3\%$ of pairs within a workload. Pooling workloads raises this to
$97.8\%$ over $498$ pairs, but part of that value comes from the separation between
workloads rather than from resolution within one, so we base the claim on the
within-workload figures. Six configurations were not reproducible across repeats, with a
ratio of largest to smallest repeated mean latency of at least $1.5$; they are retained
in the summary rather than removed, and all sit close to saturation, where neither the
analytical model nor the deployment itself is stable.
